\documentclass[10pt]{mypackage}

\usepackage[uprightgreeks,light,noDcommand]{kpfonts}
\renewcommand{\coloneq}{\coloneqq}
\renewcommand{\eqcolon}{\eqqcolon}
%\renewcommand{\mathbb}{\mathds}
%\usepackage{homework}
\usepackage{notes}

%\usepackage[ backend=bibtex, style = alphabetic, sorting=ynt ]{biblatex}
%\addbibresource{  }

\usepackage{parskip}

\fancyhf{}
\fancyhead[R]{Avinash Iyer}
\fancyhead[L]{Functional Analysis: Class Notes}
\fancyfoot[C]{\thepage}

\setcounter{secnumdepth}{0}

\begin{document}
\RaggedRight
These are notes for Malek Abdesselam's Functional Analysis class for Fall 2026.
\tableofcontents
\section{Motivation}%
Consider a space $V$ of (sufficiently nice) functions $f\colon \R\rightarrow \R$. Then, similar to how a vector in $\R^n$ can be denoted as
\begin{align*}
  x &= \begin{pmatrix}x_1\\\vdots\\x_n\end{pmatrix},
\end{align*}
we may view the function $f\colon \R\rightarrow \R$ as a ``column vector'' with uncountable index. That is, $f = \left( f(x) \right)_{x\in \R}$. If $L\colon V\rightarrow V$ is a linear map, then an analogue of the finite-dimensional matrix representation can be viewed as a kernel that $f$ is integrated against; that is,
\begin{align*}
  \left( Lf \right)(x) &= \int_{\R}^{} K\left( x,y \right)f(y)\:dy.
\end{align*}
In the special case that $L = \id_V$, then the kernel $K\left( x,y \right)$ is given by the Dirac delta distribution, $K\left( x,y \right) = \delta\left( x-y \right)$. The space of these distributions is the primary object of study in Fourier analysis.

\section{Locally Convex Topological Vector Spaces}%
Recall that if $X$ is a set, then $ \mathcal{T}\subseteq \mathcal{P}\left( \mathcal{P}\left( X \right) \right) $ is a topology on $X$ if it satisfies:
\begin{itemize}
  \item $\emptyset,X\in \mathcal{T}$;
  \item $ U_1,U_2\in \mathcal{T} $ implies $U_1\cap U_2\in \mathcal{T}$;
  \item $\set{U_i}_{i\in I}\subseteq \mathcal{T}$ implies $\bigcup_{i\in I}U_i\in \mathcal{T}$.
\end{itemize}
\begin{definition}
  We say $\left( V, \mathcal{T}\right)$ is a \textit{topological vector space} (over $\K$) if $V$ is a $\K$-vector space whose topology is such that the maps
  \begin{align*}
    V\times V &\rightarrow V\\
    \left( u,v \right) &\mapsto u + v
      \intertext{and}
    \K\times V &\rightarrow V\\
    \left( \lambda,v \right) &\mapsto \lambda v
  \end{align*}
  are continuous with respect to the product topology. Note that $\K$ is equipped with its traditional topology.
\end{definition}
The morphisms of topological vector spaces are the \textit{continuous} linear maps. We say $V\cong W$ are isomorphic as topological vector spaces if there is $\varphi\colon V\rightarrow W$ that is bijective, linear, and continuous, and $\varphi^{-1}$ is also continuous.
\begin{definition}
  Let $V$ be a $\K$-vector space. We say $\rho\colon V\rightarrow [0,\infty)$ is a \textit{seminorm} if
  \begin{align*}
    \rho\left( v + w \right) &\leq \rho\left( v \right) + \rho\left( w \right)\\
    \rho\left( \lambda v \right) &= \left\vert \lambda \right\vert \rho(v)
  \end{align*}
  for any $v,w\in V$ and $\lambda\in \K$.
\end{definition}
\begin{remark}
  The only difference between a seminorm and a norm is that $\rho(v) = 0$ need not imply that $v = 0$. A norm is a seminorm.
\end{remark}
We will let $ \operatorname{SN}(V) $ be the set of all seminorms of $V$. 
\begin{definition}
  Let $V$ be a $\K$-vector space. Let $ \mathcal{N}\subseteq \operatorname{SN}(V) $. Define $ \mathcal{T} \left( \mathcal{N} \right) $ to be a topology with basic open sets $U\left( x,F,\ve \right)\subseteq V$ given by
  \begin{align*}
    U\left( x,F,\ve \right) &= \set{y\in V : \rho\left( y-x \right) < \ve\text{ for all }\rho\in F}
  \end{align*}
  for any $x\in V$, $ F\subseteq \mathcal{N} $ finite, and $\ve > 0$.
\end{definition}
We observe that the subsets of the form $U\left( x,\set{\rho},\ve \right)$ as $\rho$ ranges across $ \mathcal{N} $ forms a subbase for the topology of $V$.
\begin{proposition}
  If $F\subseteq \mathcal{N}$ is finite, $ \ve > 0 $, and $x\in V$, then
  \begin{align*}
    U\left( x,F,\ve \right)\in \mathcal{T}\left( \mathcal{N} \right).
  \end{align*}
\end{proposition}
\begin{proof}
  Suppose $y\in U\left( x,F,\ve \right)$. Our goal is to find $\eta > 0$ such that $U\left( y,F,\eta \right) \subseteq U\left( x,F,\ve \right)$. Suppose $z\in U\left( y,F,\eta \right)$. If $\rho\in F$, then
  \begin{align*}
    \rho\left( z-x \right) &\leq \rho\left( z-y \right) + \rho\left( y-x \right)\\
                           &< \eta + \rho\left( y-x \right).
  \end{align*}
  Therefore, we will define $\eta$ to be small enough such that $\eta + \max_{\rho\in F} \rho\left( y-x \right) < \ve$. This gives the desired result.
\end{proof}
\begin{definition}
  A topological space is called a locally convex topological vector space (LCTVS) if the topology is generated by a family of seminorms $\mathcal{N}$ with the property that
  \begin{align*}
    \bigcap_{\rho\in \mathcal{N}} \set{v : \rho(v) = 0} = \set{0}.
  \end{align*}
\end{definition}
\section{Notation}%
\begin{itemize}
  \item The natural numbers are given by $\N = \set{0,1,2,3,\dots}$.
  \item The strictly positive natural numbers are given by $\Z_{> 0} = \set{1,2,3,\dots}$.
  \item A multi-index is a tuple $\alpha\in \N^{n}$. We let
    \begin{align*}
      \left\vert \alpha \right\vert = \alpha_1 + \cdots + \alpha_n.
    \end{align*}
    If $x = \left( x_1,\dots,x_n \right)\in \R^{n}$, then we will write 
    \begin{align*}
      x^{\alpha} &= x_1^{\alpha_1}\cdots x_n^{\alpha_n}
        \intertext{and}
      \alpha! &= \alpha_1!\cdots\alpha_n!
    \end{align*}
    for the corresponding values.

    If $f\colon \R^{n}\rightarrow \R$ is any function, then we will write
    \begin{align*}
      \partial^{\alpha}f &= \frac{\partial^{\left\vert \alpha \right\vert}}{\partial x_1^{\alpha_1}\cdots \partial x_n^{\alpha_n}}.
    \end{align*}
  \item If $A$ is a set, we will write $ \mathcal{P}(A) $ for the power set of $A$. If $A$ and $B$ are sets, then the set $B^{A}\subseteq \mathcal{P}\left( A\times B \right)$ is the set of all functions $f\colon A\rightarrow B$.
  \item If $n\in \N$, then
    \begin{align*}
      \left[ n \right] &= \set{1,2,\dots,n},
    \end{align*}
    where we have $\left[ 0 \right] = \emptyset$.
  \item If $A$ is a finite set, then $\left\vert A \right\vert$ is the cardinality of $A$. 
  \item Generally we will denote by $ \K $ for either $\R$ or $\C$.
\end{itemize}
\end{document}
